\section{Problems for Stage V}
\label{problems:stage-V}

\begin{enumerate}[label=\textbf{V.\arabic*},leftmargin=*]
  \item Remove the first derivative from the $d$-dimensional radial
  Schr\"odinger equation and derive the effective inverse-square coefficient.

  \item Compute the two Frobenius indices at the origin for
  $g/r^2$.  Classify the square-integrable solutions and discuss when a
  boundary condition beyond regularity is needed.

  \item Perform the Langer map $r=\rme^x$ including the wave-function
  rescaling.  Show where $(l+1/2)^2$ enters and why it is not an arbitrary
  replacement rule.

  \item Derive the three-dimensional harmonic-oscillator spectrum from a
  closed exact-WKB cycle.  Repeat with an open path and the origin monodromy.

  \item Derive the Coulomb bound spectrum from the radial action integral.
  Keep the angular-momentum and turning-point phases separate until the last
  line.

  \item Introduce a small circle $r=\epsilon\rme^{\rmi\phi}$ around the
  origin.  Calculate the local monodromy and derive a condition that makes the
  global quantization independent of $\epsilon$.

  \item Treat a weak quartic perturbation of the radial oscillator with both
  ordinary perturbation theory and the optimized monodromy scheme.  Compare
  the regimes in which each truncation is useful.

  \item Construct the turning-point equation for the Cornell potential and
  locate its roots numerically on the relevant sheet.  Test homotopy
  invariance of the quantization contour as parameters vary.

  \item Repeat the previous problem for a Yukawa-screened Coulomb potential.
  Identify the threshold at which a bound-state pole approaches the cut.

  \item Give an example in which a coordinate transformation creates an
  apparent singularity.  Distinguish its Jacobian monodromy from a physical
  self-adjoint-extension parameter.
\end{enumerate}
